% 图 1.2 GAN-based JSCC 结构图（精简专业版）
% 特点：去除杂乱装饰，避免元素重叠，保持顶会标准的简洁美学
\documentclass[tikz,border=10pt]{standalone}

\usepackage{amssymb, amsmath}
\usepackage{fontspec}
\usepackage{xeCJK}
\setCJKmainfont{SimHei}
\usepackage{tikz}
\usetikzlibrary{
  arrows.meta, positioning, shapes.geometric, calc,
  backgrounds, fit, decorations.pathreplacing
}

% 低饱和度马卡龙配色
\definecolor{mBlue}{HTML}{AEC6CF}
\definecolor{mPink}{HTML}{FFD1DC}
\definecolor{mGreen}{HTML}{B1D8B7}
\definecolor{mPurple}{HTML}{B39EB5}
\definecolor{mYellow}{HTML}{FDFD96}
\definecolor{mGray}{HTML}{E5E4E2}
\definecolor{mDark}{HTML}{4A4A4A}

\tikzset{
  >=stealth, thick, draw=mDark,
  % 主要模块：统一大小，避免视觉混乱
  mainblock/.style={
    rectangle, rounded corners=3pt, draw=mDark, align=center,
    font=\sffamily\small, minimum width=2.2cm, minimum height=0.9cm, inner sep=5pt},
  % 微观模块：简化，不完全展开
  micro/.style={
    rectangle, rounded corners=2pt, draw=mDark!70, align=center,
    font=\sffamily\tiny, minimum height=0.4cm, inner sep=2pt},
  % 张量标注：统一样式
  ts/.style={font=\tiny\sffamily, color=mDark!75},
  % 箭头：统一粗细
  arr/.style={-{Stealth[length=4pt]}, thick, draw=mDark},
  darr/.style={-{Stealth[length=4pt]}, thick, draw=mDark, dashed},
  % 图例：简洁
  legendbox/.style={
    rectangle, draw=mDark!50, rounded corners=3pt,
    fill=mGray!15, inner sep=5pt, font=\tiny\sffamily},
}

\begin{document}
\begin{tikzpicture}[
  node distance=0.7cm and 1.3cm,
]
  % ========== 主数据流（单行布局，避免重叠）==========
  \node[mainblock, fill=mGray!50] (img) {$\mathbf{x}$};
  \node[mainblock, fill=mBlue!40, right=1.3cm of img] (enc) {Encoder};
  \node[mainblock, fill=mBlue!40, right=1.3cm of enc] (pwr) {功率控制};
  \node[mainblock, fill=mGreen!50, right=1.3cm of pwr] (awgn) {AWGN};
  \node[mainblock, fill=mPurple!40, right=1.3cm of awgn] (gen) {Generator};
  \node[mainblock, fill=mGray!50, right=1.3cm of gen] (out) {$\hat{\mathbf{x}}$};

  % ========== 主流程连线（单层标注，避免拥挤）==========
  \draw[arr] (img) -- (enc)
    node[pos=0.5, above=3pt, ts] {$1{\times}3{\times}256{\times}256$};

  \draw[arr] (enc) -- (pwr)
    node[pos=0.5, above=3pt, ts] {$\mathbf{y} \in \mathbb{R}^{1{\times}32^3}$};

  \draw[arr] (pwr) -- (awgn)
    node[pos=0.5, above=3pt, ts] {$\sqrt{P}\mathbf{y}$};

  \draw[arr] (awgn) -- (gen)
    node[pos=0.5, above=3pt, ts] {$\hat{\mathbf{y}} = \sqrt{P}\mathbf{y} + \mathbf{n}$};

  \draw[arr] (gen) -- (out)
    node[pos=0.5, above=3pt, ts] {$1{\times}3{\times}256{\times}256$};

  % ========== Discriminator（简化布局）==========
  \node[mainblock, fill=mPink!40, below=1.8cm of awgn, dashed, minimum width=2.5cm] (disc)
    {Discriminator\\{\tiny 仅训练}};

  \draw[darr] (img.south) -- ++(0,-0.5) -| (disc.west)
    node[pos=0.15, left=2pt, ts] {$\mathbf{x}$};

  \draw[darr] (out.south) -- ++(0,-0.5) -| (disc.east)
    node[pos=0.15, right=2pt, ts] {$\hat{\mathbf{x}}$};

  \draw[darr] (awgn.south) -- (disc.north)
    node[pos=0.5, right=2pt, ts] {$\hat{\mathbf{y}}$};

  % ========== Encoder 微观结构（精简，不完全展开）==========
  \node[micro, fill=mBlue!30, below=1.0cm of enc, xshift=-1.0cm] (e1) {Conv};
  \node[micro, fill=mBlue!30, right=0.15cm of e1] (e2) {BN};
  \node[micro, fill=mBlue!30, right=0.15cm of e2] (e3) {ReLU};
  \node[micro, fill=mBlue!25, right=0.2cm of e3, minimum width=1.0cm] (e4) {Res$\times$2};

  \draw[arr, -, thin] (e1) -- (e2) -- (e3) -- (e4);
  \draw[arr, dashed, thin, mDark!60] (enc.south) -- (e2.north);

  \node[ts, below=0.2cm of e1] {$3{\to}64$};
  \node[ts, below=0.2cm of e4] {$256{\to}32$};

  % ========== Generator 微观结构（精简）==========
  \node[micro, fill=mPurple!30, below=1.0cm of gen, xshift=-1.0cm] (g1) {Conv};
  \node[micro, fill=mPurple!30, right=0.15cm of g1] (g2) {BN};
  \node[micro, fill=mPurple!30, right=0.15cm of g2] (g3) {ReLU};
  \node[micro, fill=mPurple!25, right=0.2cm of g3, minimum width=1.0cm] (g4) {Res$\times$5};

  \draw[arr, -, thin] (g1) -- (g2) -- (g3) -- (g4);
  \draw[arr, dashed, thin, mDark!60] (gen.south) -- (g2.north);

  \node[ts, below=0.2cm of g1] {$32{\to}256$};
  \node[ts, below=0.2cm of g4] {$256{\to}3$};

  % ========== 背景分组（淡化，避免视觉干扰）==========
  \begin{scope}[on background layer]
    \node[fit=(img)(enc)(pwr)(e1)(e4),
      fill=mBlue!8, draw=mDark!30, dashed,
      rounded corners=5pt, inner sep=8pt] {};

    \node[fit=(gen)(out)(g1)(g4),
      fill=mPurple!8, draw=mDark!30, dashed,
      rounded corners=5pt, inner sep=8pt] {};
  \end{scope}

  % 标签（简洁，不遮挡）
  \node[above=0.35cm of enc, font=\sffamily\scriptsize, color=mDark!80] {发送端};
  \node[above=0.35cm of gen, font=\sffamily\scriptsize, color=mDark!80] {接收端};

  % ========== 核心数学公式（左上角，紧凑排版）==========
  \node[legendbox, anchor=north west] at ([xshift=5pt, yshift=-5pt]current bounding box.north west)
    [align=left, text width=4.2cm] {
    \textbf{训练损失}\\[2pt]
    $L_t = \lambda L_G + \alpha L_{\text{MSE}} + \beta L_{\text{LPIPS}}$\\[1pt]
    $L_G = -\mathbb{E}[\log D(G(\hat{\mathbf{y}}), \hat{\mathbf{y}})]$\\[1pt]
    $L_D = -\mathbb{E}[\log D(\mathbf{x}, \mathbf{y})]$\\
    \hspace{1.2cm}$- \mathbb{E}[\log(1{-}D(\hat{\mathbf{x}}, \hat{\mathbf{y}}))]$
  };

  % ========== 图例（右下角，精简）==========
  \node[legendbox, anchor=south east] at ([xshift=-5pt, yshift=5pt]current bounding box.south east)
    [align=left, text width=4.0cm] {
    \textbf{图例}\\[1pt]
    \colorbox{mBlue!40}{\rule{0.15cm}{0.1cm}} 编码 \quad
    \colorbox{mPurple!40}{\rule{0.15cm}{0.1cm}} 解码 \quad
    \colorbox{mGreen!50}{\rule{0.15cm}{0.1cm}} 信道\\[1pt]
    实线：推理 \quad 虚线：训练\\[2pt]
    SNR $= 10\log_{10}(\|\mathbf{y}\|^2/\sigma^2)$ dB
  };
\end{tikzpicture}
\end{document}
